\section{Modelling transformations}
\label{sec:trs}
\here{2}
\concept{Each part's local matrix places a unit primitive: scale first, then rotate, then translate. In homogeneous
coordinates all three are $4\times4$ matrices, so one product gives the whole placement.}

\begin{gather}
M = T(\vect t)\,R_y(\psi)\,R_x(\theta)\,R_z(\varphi)\,S(\vect s)\label{eq:trs}\\
T=\begin{bmatrix}1&0&0&t_x\\0&1&0&t_y\\0&0&1&t_z\\0&0&0&1\end{bmatrix},\
R_y=\begin{bmatrix}\cos\psi&0&\sin\psi&0\\0&1&0&0\\-\sin\psi&0&\cos\psi&0\\0&0&0&1\end{bmatrix},\
S=\mathrm{diag}(s_x,s_y,s_z,1)\nonumber
\end{gather}
\begin{equation}
\text{rotation about a pivot } \vect p:\quad M_{\text{pivot}} = T(\vect p)\,R\,T(-\vect p)
\label{eq:pivot}
\end{equation}
Order matters: $T R S \neq S R T$. Scaling after a rotation would shear a rotated box; translating before a rotation
would swing the part around the parent's origin. The turret, gun barrel, rotor and door hinges use
Eq.~\eqref{eq:pivot}: each sits in a group node whose origin is the pivot.

\begin{lstlisting}[language=C++,caption={The local matrix of every scene node (\code{src/scene/Transform.h}).}]
// Model matrix M = T * R * S with R = Ry * Rx * Rz: the primitive is scaled first,
// then rotated (roll, pitch, yaw), then translated.
Mat4 matrix() const { return translate(translation) * rotateEulerYXZ(rotationDeg) * ts::scale(scale); }
\end{lstlisting}

\begin{figure}[H]
\centering
\raw{scene/soldier_axes}{0.72\linewidth}
\caption{Local versus world coordinates (\key{H}: local axes of the hierarchy), shown on a soldier's joints: each
red/green/blue triad is one local frame. Notice that the axis triads are rotated with their parts: each part's
geometry is defined in its own frame, and only the product of matrices up the hierarchy places it in the world.}
\label{fig:axes}
\end{figure}

\begin{figure}[H]
\centering
{\footnotesize\textit{Inputs exported by the program (part UAV-14, parent UAV-00)}}\\[0.5mm]
\wbblock{Transform_Trace}{inputs}\\[1mm]
{\footnotesize\textit{Scale and the three rotations}}\\[0.5mm]
\wbblock{Transform_Trace}{srt}
{\footnotesize\textit{Translation, local matrix, parent and world matrices}}\\[0.5mm]
\wbblock{Transform_Trace}{tm}\wbblock{Transform_Trace}{tw}
\begin{minipage}[t]{0.6\linewidth}\centering{\footnotesize\textit{The vertex through the pipeline}}\\[0.5mm]
\wbblock{Transform_Trace}{stages}\end{minipage}\hfill
\begin{minipage}[t]{0.38\linewidth}\centering{\footnotesize\textit{Checks}}\\[0.5mm]
\wbblock{Transform_Trace}{checks}\end{minipage}\\[1mm]
{\footnotesize\wbkey}
\caption{Workbook sheet \code{Transform\_Trace}: one real vertex of one real UAV part traced through $S$, $R$, $T$,
the parent matrix, view, projection, perspective divide, NDC and viewport, each matrix as a live $4\times4$ grid.
Notice the check column: the world matrix computed in the sheet matches the one the program exports. \fullsheet{Transform_Trace}}
\label{fig:wb-trace}
\end{figure}

\where{src/scene/Transform.h (matrix), src/math/Math3D.cpp (translate, rotateEulerYXZ, scale), src/objects/Builder.cpp}
\critique{Euler angles in a fixed order are easy to tabulate (the brief's T/R/S table) but suffer gimbal lock at
$\theta=\pm90^\circ$; the camera rig and the interpolation between simulation ticks therefore use quaternions
(\code{RigPose.rot}). Reflection and shear are not used by any part.}
